% Generated draft: requires mathematical and source review.
\documentclass[11pt]{article}
\usepackage[margin=25mm]{geometry}
\usepackage{fontspec}
\setmainfont{Latin Modern Roman}
\usepackage{amsmath,amssymb,mathtools,unicode-math}
\setmathfont{Latin Modern Math}
\usepackage{enumitem,needspace,xurl,hyperref}
\setlength{\parindent}{0pt}
\setlength{\parskip}{5pt}
\setlength{\emergencystretch}{3em}
\newcommand{\N}{\mathbb{N}}
\newcommand{\Z}{\mathbb{Z}}
\newcommand{\Q}{\mathbb{Q}}
\newcommand{\R}{\mathbb{R}}
\newcommand{\C}{\mathbb{C}}
\newcommand{\F}{\mathbb{F}}
\newcommand{\A}{\mathbb{A}}
\newcommand{\PP}{\mathbb{P}}
\newcommand{\E}{\mathbb{E}}
\newcommand{\eps}{\varepsilon}
\newcommand{\dd}{\,\mathrm{d}}
\begin{document}
\section*{UnsolvedMath ID: 30005352}
\textbf{Source problem code:} catalogue entry
\subsection{Definitions and mathematical statement}
Let \(\F_2\) be the field with two elements. For each positive integer \(n\), a two-coloring is any map \(c:\F_2^n\setminus\{0\}\to\{1,2\}\); both colors need not occur. A nonzero linear subspace \(V\leq\F_2^n\) is monochromatic away from zero if there exists \(i\in\{1,2\}\) such that \(c(v)=i\) for every \(v\in V\setminus\{0\}\). Define
\[
M(c)=\max\bigl\{|V|-1: \{0\}\ne V\leq\F_2^n,\ V\setminus\{0\}\text{ is monochromatic under }c\bigr\},
\qquad
f(n)=\min_{c:\F_2^n\setminus\{0\}\to\{1,2\}}M(c).
\]
These extrema exist: there are finitely many colorings and subspaces, and every one-dimensional subspace has a monochromatic singleton of nonzero vectors. A subspace of dimension \(d\) contributes \(2^d-1\) nonzero vectors. Thus \(f(n)\) counts nonzero vectors, rather than dimension. This definition also agrees with the supplied largest guaranteed exact cardinality: each \(f(n)\) has the form \(2^d-1\), and any larger monochromatic subspace contains a subspace of dimension \(d\).

The question is whether
\[
\exists C>0\ \exists n_0\in\Z_{\geq1}\ \forall n\in\Z_{\geq1},\qquad n\geq n_0\Longrightarrow f(n)\leq Cn.
\]
Equivalently, does there exist such a constant \(C\) so that, for every sufficiently large positive integer \(n\), some coloring of \(\F_2^n\setminus\{0\}\) has at most \(Cn\) nonzero vectors in every monochromatic linear subspace? The zero vector is excluded from coloring and from the cardinality being bounded. The supplied record does not specify an \(n=0\) convention; none is needed for this asymptotic question.

There is an unresolved source-extraction qualification. The supplied clean statement and title identify the question above, but the raw original statement contains only a garbled terminal question about \(f(n)=O(n)\), followed by a different Question 12 headed Karim Adiprasito. That graph question concerns attaching disks along a collection \(\mathcal C\) of cycles in a graph \(G\), minimizing the total cycle length \(\sum_{\gamma\in\mathcal C}|\gamma|\) to kill homology over a field, integral homology, the first homotopy group, or all homotopy groups. Its displayed hierarchy is \(\alpha_F(G)\leq\alpha_{\Z}(G)\leq\alpha_{\pi_1}(G)\leq\alpha_{\infty}(G)\), and it asks whether \(\alpha_{\pi_1}(G)/\alpha_{\F_2}(G)\) is unbounded, with bounded degree requested in the strongest form. This neighboring material does not supply the missing raw definition of \(f\), and its attribution cannot safely be transferred to the vector-space question. The formulation above therefore follows the supplied clean statement, subject to human source review.
\subsection{Short English statement}
For every positive integer \(n\), let \(f(n)\) be the largest number of nonzero vectors guaranteed in a monochromatic linear subspace under every two-coloring of \(\F_2^n\setminus\{0\}\). Is \(f(n)\) bounded above by a constant times \(n\) for all sufficiently large \(n\)? The count concerns vectors, not dimension. The supplied raw extraction includes a neighboring graph problem, so the clean statement's source correspondence requires review.
\subsection{Sources}
\begin{itemize}
\item Supplied local UnsolvedMath registry record, permanent ID \(30005352\), title \emph{Monochromatic Subspaces in Binary Vector Spaces}, exact source problem code \texttt{catalogue entry}. The record gives proposed year 2023 and no proposer for this problem.
\item \emph{Combinatorics (2023)}, Oberwolfach Reports, pp.\ 5--89, DOI 10.4171/owr/2023/1: \url{https://doi.org/10.4171/owr/2023/1}. The supplied background describes Oberwolfach Reports as documenting workshops and their open-problem discussions. No precise page for this individual question is supplied.
\item Conlon--Fox--Pham--Yepremyan, \emph{Random Cayley Graphs}: \url{https://arxiv.org/abs/2412.21194}. This reference is supplied as literature context; its exact relevant result is not specified in the record.
\item Supplied dated literature curation, checked on \texttt{2026-08-22}: status \texttt{open}. Its assessment says that the specified binary-vector-space Ramsey problem still has only an upper bound with a logarithmic factor in the cited record, and that recent results for almost all groups do not establish \(f(n)=O(n)\) for every \(n\). This is supplied dated curation derived from the OWR record, not an independently verified literature assessment; the background explicitly requests independent verification of time-sensitive status and cited claims.
\item The registry labels the statement \texttt{corrected\_verified} and says it was checked against the official Oberwolfach report, preserving its hypotheses, notation, and open target. These are supplied verification assertions, not verification performed for this draft. The raw extraction's neighboring attribution to Karim Adiprasito remains separate from the vector-space problem.
\end{itemize}
\end{document}
